% GENERATED FILE. Edit canonical lessons, then run npm run generate:books.
% canonical-source-hash: fnv1a64:22828b3f74ee3a67
% canonical-lessons: TA-C500-tunai, TA-C500-illattaraci, TA-C500-pukunta-vitu, TA-C500-tanikkutittanam, TA-C500-ottuppo

\chapter{Running a Home}
\label{ch:ta-running-a-home}
\hlchaptermodality{500}

\begin{chapteropening}
\textbf{By the end of this chapter:} \emph{I can talk about companions, homemakers, the marital home and setting up house.}

Complete the last lesson of chapter 500: I can talk about companions, homemakers, the marital home and setting up house.
\end{chapteropening}

\section[\texorpdfstring{\ta{துணை} (tuṇai)}{துணை (tuṇai)}]{\ta{துணை} (tuṇai) — a companion; a partner}

\label{lesson:TA-C500-tunai}



\begin{quote}
Before the new one: say the Tamil for a married couple, then the Tamil for married life.
\end{quote}

\subsection*{You'll want to know: \ta{துணை}}
\textbf{\ta{துணை}} — \emph{tuṇai} — \textquotedblleft{}a companion; a partner\textquotedblright{}.

\begin{quote}
\ta{பயணத்தில்} \ta{எனக்குத்} \ta{துணையாக} \ta{என்} \ta{தம்பி} \ta{வந்தான்}.
\end{quote}

\begin{quote}
\emph{payaṇattil eṉakkut tuṇaiyāka eṉ tampi vantāṉ.}
\end{quote}

\begin{quote}
My younger brother came along to keep me company on the journey.
\end{quote}

\ta{துணையாக}, as company; \ta{வாழ்க்கைத்} \ta{துணை}, a life partner.

\begin{cousinweb}
\textbf{\ta{துணை}} (\emph{tuṇai}) is native Dravidian (DEDR 3308), for company, help and a companion.
\end{cousinweb}

\begin{grammarlens}[title={What you've built}]
One more word for describing what you lived and felt.
\end{grammarlens}

\subsection*{Your turn}
\begin{itemize}
\raggedright
  \item \emph{Say it:} \emph{tuṇai}
  \item \emph{Say it:} \emph{tuṇai}, once more
  \item \emph{Say it:} \emph{payaṇattil eṉakkut tuṇaiyāka eṉ tampi vantāṉ.} — My younger brother came along to keep me company on the journey.
  \item \emph{Recall:} say the Tamil for a married couple, then the Tamil for married life, then say \emph{tuṇai} again
\end{itemize}

\begin{tcolorbox}[breakable,colback=teal!4,colframe=teal!35!black,title={Before you move on}]
What does \ta{துணை} mean? (\textquotedblleft{}A companion; a partner\textquotedblright{}.) Say it once more.
\end{tcolorbox}
\section[\texorpdfstring{\ta{இல்லத்தரசி} (illattaraci)}{இல்லத்தரசி (illattaraci)}]{\ta{இல்லத்தரசி} (illattaraci) — a homemaker; a housewife}

\label{lesson:TA-C500-illattaraci}



\begin{quote}
Before the new one: say the Tamil for married life, then the Tamil for a companion.
\end{quote}

\subsection*{You'll want to know: \ta{இல்லத்தரசி}}
\textbf{\ta{இல்லத்தரசி}} — \emph{illattaraci} — \textquotedblleft{}a homemaker; a housewife\textquotedblright{}.

\begin{quote}
\ta{என்} \ta{அம்மா} \ta{ஒரு} \ta{இல்லத்தரசி}.
\end{quote}

\begin{quote}
\emph{eṉ ammā oru illattaraci.}
\end{quote}

\begin{quote}
My mother is a homemaker.
\end{quote}

\ta{இல்லம்} \textquotedblleft{}home\textquotedblright{} + \ta{அரசி} \textquotedblleft{}queen\textquotedblright{}.

\begin{cousinweb}
\textbf{\ta{இல்லத்தரசி}} (\emph{illattaraci}) is a compound of native \ta{இல்லம்} \emph{illam}, \textquotedblleft{}home\textquotedblright{} (DEDR 494), and \ta{அரசி} \emph{araci}, \textquotedblleft{}queen\textquotedblright{}, the feminine of \ta{அரசன்} \textquotedblleft{}king\textquotedblright{}, which DEDR 201 traces to Sanskrit \emph{rājan}.
\end{cousinweb}

\begin{grammarlens}[title={What you've built}]
One more word for describing what you lived and felt.
\end{grammarlens}

\subsection*{Your turn}
\begin{itemize}
\raggedright
  \item \emph{Say it:} \emph{illattaraci}
  \item \emph{Say it:} \emph{illattaraci}, once more
  \item \emph{Say it:} \emph{eṉ ammā oru illattaraci.} — My mother is a homemaker.
  \item \emph{Recall:} say the Tamil for married life, then the Tamil for a companion, then say \emph{illattaraci} again
\end{itemize}

\begin{tcolorbox}[breakable,colback=teal!4,colframe=teal!35!black,title={Before you move on}]
What does \ta{இல்லத்தரசி} mean? (\textquotedblleft{}A homemaker; a housewife\textquotedblright{}.) Say it once more.
\end{tcolorbox}
\section[\texorpdfstring{\ta{புகுந்த} \ta{வீடு} (pukunta vīṭu)}{புகுந்த வீடு (pukunta vīṭu)}]{\ta{புகுந்த} \ta{வீடு} (pukunta vīṭu) — a woman’s marital home}

\label{lesson:TA-C500-pukunta-vitu}



\begin{quote}
Before the new one: say the Tamil for a companion, then the Tamil for a homemaker.
\end{quote}

\subsection*{You'll want to know: \ta{புகுந்த} \ta{வீடு}}
\textbf{\ta{புகுந்த} \ta{வீடு}} — \emph{pukunta vīṭu} — \textquotedblleft{}a woman’s marital home\textquotedblright{}.

\begin{quote}
\ta{திருமணத்துக்குப்} \ta{பிறகு} \ta{அவள்} \ta{புகுந்த} \ta{வீட்டுக்குப்} \ta{போனாள்}.
\end{quote}

\begin{quote}
\emph{tirumaṇattukkup piṟaku avaḷ pukunta vīṭṭukkup pōṉāḷ.}
\end{quote}

\begin{quote}
After the wedding she went to her husband’s home.
\end{quote}

\ta{புகு} \textquotedblleft{}to enter\textquotedblright{}: the house she entered.

\begin{cousinweb}
\textbf{\ta{புகுந்த} \ta{வீடு}} (\emph{pukunta vīṭu}) is native: \ta{புகுந்த}, \textquotedblleft{}entered\textquotedblright{}, from \ta{புகு} \emph{puku}, \textquotedblleft{}to enter\textquotedblright{} (DEDR 4238), and \ta{வீடு}, \textquotedblleft{}house\textquotedblright{}.
\end{cousinweb}

\begin{grammarlens}[title={What you've built}]
One more word for describing what you lived and felt.
\end{grammarlens}

\subsection*{Your turn}
\begin{itemize}
\raggedright
  \item \emph{Say it:} \emph{pukunta vīṭu}
  \item \emph{Say it:} \emph{pukunta vīṭu}, once more
  \item \emph{Say it:} \emph{tirumaṇattukkup piṟaku avaḷ pukunta vīṭṭukkup pōṉāḷ.} — After the wedding she went to her husband’s home.
  \item \emph{Recall:} say the Tamil for a companion, then the Tamil for a homemaker, then say \emph{pukunta vīṭu} again
\end{itemize}

\begin{tcolorbox}[breakable,colback=teal!4,colframe=teal!35!black,title={Before you move on}]
What does \ta{புகுந்த} \ta{வீடு} mean? (\textquotedblleft{}A woman’s marital home\textquotedblright{}.) Say it once more.
\end{tcolorbox}
\section[\texorpdfstring{\ta{தனிக்குடித்தனம்} (taṉikkuṭittaṉam)}{தனிக்குடித்தனம் (taṉikkuṭittaṉam)}]{\ta{தனிக்குடித்தனம்} (taṉikkuṭittaṉam) — a separate household}

\label{lesson:TA-C500-tanikkutittanam}



\begin{quote}
Before the new one: say the Tamil for a homemaker, then the Tamil for a woman’s marital home.
\end{quote}

\subsection*{You'll want to know: \ta{தனிக்குடித்தனம்}}
\textbf{\ta{தனிக்குடித்தனம்}} — \emph{taṉikkuṭittaṉam} — \textquotedblleft{}a separate household\textquotedblright{}.

\begin{quote}
\ta{திருமணத்துக்குப்} \ta{பிறகு} \ta{அவர்கள்} \ta{தனிக்குடித்தனம்} \ta{போனார்கள்}.
\end{quote}

\begin{quote}
\emph{tirumaṇattukkup piṟaku avarkaḷ taṉikkuṭittaṉam pōṉārkaḷ.}
\end{quote}

\begin{quote}
After the wedding they set up a home of their own.
\end{quote}

\ta{தனி} \textquotedblleft{}separate\textquotedblright{} + \ta{குடித்தனம்} \textquotedblleft{}household\textquotedblright{}; the opposite of living as a joint family.

\begin{cousinweb}
\textbf{\ta{தனிக்குடித்தனம்}} (\emph{taṉikkuṭittaṉam}) is a native compound of \ta{தனி} \emph{taṉi}, \textquotedblleft{}alone, separate\textquotedblright{} (DEDR 3196), and \ta{குடித்தனம்}, \textquotedblleft{}a household\textquotedblright{}.
\end{cousinweb}

\begin{grammarlens}[title={What you've built}]
One more word for describing what you lived and felt.
\end{grammarlens}

\subsection*{Your turn}
\begin{itemize}
\raggedright
  \item \emph{Say it:} \emph{taṉikkuṭittaṉam}
  \item \emph{Say it:} \emph{taṉikkuṭittaṉam}, once more
  \item \emph{Say it:} \emph{tirumaṇattukkup piṟaku avarkaḷ taṉikkuṭittaṉam pōṉārkaḷ.} — After the wedding they set up a home of their own.
  \item \emph{Recall:} say the Tamil for a homemaker, then the Tamil for a woman’s marital home, then say \emph{taṉikkuṭittaṉam} again
\end{itemize}

\begin{tcolorbox}[breakable,colback=teal!4,colframe=teal!35!black,title={Before you move on}]
What does \ta{தனிக்குடித்தனம்} mean? (\textquotedblleft{}A separate household\textquotedblright{}.) Say it once more.
\end{tcolorbox}
\section[\texorpdfstring{\ta{ஒத்துப்போ} (ottuppō)}{ஒத்துப்போ (ottuppō)}]{\ta{ஒத்துப்போ} (ottuppō) — to get along}

\label{lesson:TA-C500-ottuppo}



\begin{quote}
Before the new one: say the Tamil for a woman’s marital home, then the Tamil for a separate household.
\end{quote}

\subsection*{You'll want to know: \ta{ஒத்துப்போ}}
\textbf{\ta{ஒத்துப்போ}} — \emph{ottuppō} — \textquotedblleft{}to get along\textquotedblright{}.

\begin{quote}
\ta{என்} \ta{அம்மாவும்} \ta{பாட்டியும்} \ta{நன்றாக} \ta{ஒத்துப்போனார்கள்}.
\end{quote}

\begin{quote}
\emph{eṉ ammāvum pāṭṭiyum naṉṟāka ottuppōṉārkaḷ.}
\end{quote}

\begin{quote}
My mother and grandmother got along well.
\end{quote}

Past (I …): \ta{ஒத்துப்போனேன்}, \emph{ottuppōṉēṉ}.

With -\ta{உடன்}: \ta{அவருடன்} \ta{ஒத்துப்போ}, to get on with him.

\begin{cousinweb}
\textbf{\ta{ஒத்துப்போ}} (\emph{ottuppō}) is built from native words: \ta{ஒத்து}, \textquotedblleft{}agreeing\textquotedblright{}, the participle of \ta{ஒ} \emph{o}, \textquotedblleft{}to suit, to be in concord\textquotedblright{} (DEDR 924), and \ta{போ}, \textquotedblleft{}to go\textquotedblright{}.
\end{cousinweb}

\begin{grammarlens}[title={What you've built}]
One more word for describing what you lived and felt.
\end{grammarlens}

\subsection*{Your turn}
\begin{itemize}
\raggedright
  \item \emph{Say it:} \emph{ottuppō}
  \item \emph{Say it:} \emph{ottuppō}, once more
  \item \emph{Say it:} \emph{eṉ ammāvum pāṭṭiyum naṉṟāka ottuppōṉārkaḷ.} — My mother and grandmother got along well.
  \item \emph{Recall:} say the Tamil for a woman’s marital home, then the Tamil for a separate household, then say \emph{ottuppō} again
\end{itemize}

\begin{tcolorbox}[breakable,colback=teal!4,colframe=teal!35!black,title={Before you move on}]
What does \ta{ஒத்துப்போ} mean? (\textquotedblleft{}To get along\textquotedblright{}.) Say it once more.
\end{tcolorbox}
